% GENERATED FILE. Edit canonical lessons, then run npm run generate:books.
% canonical-source-hash: fnv1a64:bb8fb40dff8231c1
% canonical-lessons: IT-C535-trapano, IT-C535-vite, IT-C535-cacciavite, IT-C535-aggiustare, IT-C535-imbiancare

\chapter{Fai da te}
\label{ch:it-fai-da-te}
\hlchaptermodality{535}

\begin{chapteropening}
\textbf{By the end of this chapter:} \emph{I can talk about small repairs and painting the walls.}

Complete the last lesson of chapter 535: I can talk about small repairs and painting the walls.
\end{chapteropening}

\section[il trapano]{il trapano — drill}

\label{lesson:IT-C535-trapano}



\begin{quote}
Before the new one: say the Italian for sink, then the Italian for to polish.
\end{quote}

\subsection*{You'll want to know: il trapano}
\textbf{il trapano} — \textquotedblleft{}drill\textquotedblright{}.

\begin{quote}
Mi presti il trapano? Devo fare un buco nel muro.
\end{quote}

\begin{quote}
\emph{Can you lend me the drill? I have to make a hole in the wall.}
\end{quote}

\emph{Il trapano elettrico} — the power drill.

\begin{cousinweb}
\emph{trapano} comes from Latin \emph{trepanum} \textquotedblleft{}a drill\textquotedblright{}, from Greek \emph{trýpanon} \textquotedblleft{}an auger, a borer\textquotedblright{} (EtymDB, under English \emph{trepan}).
\end{cousinweb}

\begin{grammarlens}[title={What you've built}]
One more word for telling, explaining and describing at B1.
\end{grammarlens}

\subsection*{Your turn}
\begin{itemize}
\raggedright
  \item \emph{Say it:} \emph{il trapano}
  \item \emph{Say it:} \emph{il trapano}, once more
  \item \emph{Say it:} a sentence of your own with \emph{il trapano}
  \item \emph{Recall:} say the Italian for sink, then the Italian for to polish, then say \emph{il trapano} again
\end{itemize}

\begin{tcolorbox}[breakable,colback=teal!4,colframe=teal!35!black,title={Before you move on}]
What does il trapano mean? (\textquotedblleft{}Drill\textquotedblright{}.) Say it once more.
\end{tcolorbox}
\section[la vite]{la vite — screw; vine}

\label{lesson:IT-C535-vite}



\begin{quote}
Before the new one: say the Italian for to polish, then the Italian for drill.
\end{quote}

\subsection*{You'll want to know: la vite}
\textbf{la vite} — \textquotedblleft{}screw; vine\textquotedblright{}.

\begin{quote}
Questa vite è troppo lunga.
\end{quote}

\begin{quote}
\emph{This screw is too long.}
\end{quote}

Also the grapevine: \emph{la vite} gives grapes.

\begin{cousinweb}
\emph{vite} is inherited from Latin \emph{vītis} \textquotedblleft{}a vine\textquotedblright{} (EtymDB): the screw is probably named for the spiral of the vine's tendrils.
\end{cousinweb}

\begin{grammarlens}[title={What you've built}]
One more word for telling, explaining and describing at B1.
\end{grammarlens}

\subsection*{Your turn}
\begin{itemize}
\raggedright
  \item \emph{Say it:} \emph{la vite}
  \item \emph{Say it:} \emph{la vite}, once more
  \item \emph{Say it:} a sentence of your own with \emph{la vite}
  \item \emph{Recall:} say the Italian for to polish, then the Italian for drill, then say \emph{la vite} again
\end{itemize}

\begin{tcolorbox}[breakable,colback=teal!4,colframe=teal!35!black,title={Before you move on}]
What does la vite mean? (\textquotedblleft{}Screw; vine\textquotedblright{}.) Say it once more.
\end{tcolorbox}
\section[il cacciavite]{il cacciavite — screwdriver}

\label{lesson:IT-C535-cacciavite}



\begin{quote}
Before the new one: say the Italian for drill, then the Italian for screw.
\end{quote}

\subsection*{You'll want to know: il cacciavite}
\textbf{il cacciavite} — \textquotedblleft{}screwdriver\textquotedblright{}.

\begin{quote}
Mi serve un cacciavite piccolo.
\end{quote}

\begin{quote}
\emph{I need a small screwdriver.}
\end{quote}

Literally \textquotedblleft{}drive-screw\textquotedblright{}; the same form in the plural.

\begin{cousinweb}
\emph{cacciavite} is built from \emph{caccia} \textquotedblleft{}drives\textquotedblright{}, from \emph{cacciare}, and \emph{vite} \textquotedblleft{}screw\textquotedblright{}, inherited from Latin \emph{vītis} \textquotedblleft{}a vine\textquotedblright{} (EtymDB).
\end{cousinweb}

\begin{grammarlens}[title={What you've built}]
One more word for telling, explaining and describing at B1.
\end{grammarlens}

\subsection*{Your turn}
\begin{itemize}
\raggedright
  \item \emph{Say it:} \emph{il cacciavite}
  \item \emph{Say it:} \emph{il cacciavite}, once more
  \item \emph{Say it:} a sentence of your own with \emph{il cacciavite}
  \item \emph{Recall:} say the Italian for drill, then the Italian for screw, then say \emph{il cacciavite} again
\end{itemize}

\begin{tcolorbox}[breakable,colback=teal!4,colframe=teal!35!black,title={Before you move on}]
What does il cacciavite mean? (\textquotedblleft{}Screwdriver\textquotedblright{}.) Say it once more.
\end{tcolorbox}
\section[aggiustare]{aggiustare — to fix, to mend}

\label{lesson:IT-C535-aggiustare}



\begin{quote}
Before the new one: say the Italian for screw, then the Italian for screwdriver.
\end{quote}

\subsection*{You'll want to know: aggiustare}
\textbf{aggiustare} — \textquotedblleft{}to fix, to mend\textquotedblright{}.

\begin{quote}
Mio padre aggiusta tutto in casa.
\end{quote}

\begin{quote}
\emph{My father fixes everything in the house.}
\end{quote}

Close to \emph{riparare}; also to adjust.

\begin{cousinweb}
\emph{aggiustare} is built on \emph{giusto} \textquotedblleft{}right, exact\textquotedblright{}, inherited from Latin \emph{iūstus} \textquotedblleft{}right, lawful\textquotedblright{} (Lewis \& Short): to make right.
\end{cousinweb}

\begin{grammarlens}[title={What you've built}]
One more word for telling, explaining and describing at B1.
\end{grammarlens}

\subsection*{Your turn}
\begin{itemize}
\raggedright
  \item \emph{Say it:} \emph{aggiustare}
  \item \emph{Say it:} \emph{aggiustare}, once more
  \item \emph{Say it:} a sentence of your own with \emph{aggiustare}
  \item \emph{Recall:} say the Italian for screw, then the Italian for screwdriver, then say \emph{aggiustare} again
\end{itemize}

\begin{tcolorbox}[breakable,colback=teal!4,colframe=teal!35!black,title={Before you move on}]
What does aggiustare mean? (\textquotedblleft{}To fix, to mend\textquotedblright{}.) Say it once more.
\end{tcolorbox}
\section[imbiancare]{imbiancare — to paint (walls), to whitewash}

\label{lesson:IT-C535-imbiancare}



\begin{quote}
Before the new one: say the Italian for screwdriver, then the Italian for to fix.
\end{quote}

\subsection*{You'll want to know: imbiancare}
\textbf{imbiancare} — \textquotedblleft{}to paint (walls), to whitewash\textquotedblright{}.

\begin{quote}
Questo fine settimana imbianchiamo la camera.
\end{quote}

\begin{quote}
\emph{This weekend we are painting the bedroom.}
\end{quote}

Built on \emph{bianco} \textquotedblleft{}white\textquotedblright{}, even if the paint is another colour.

\begin{cousinweb}
\emph{imbiancare} is built on \emph{bianco} \textquotedblleft{}white\textquotedblright{}, which comes from Proto-Germanic \emph{blankaz} \textquotedblleft{}bright, shining, white\textquotedblright{} (EtymDB, under Spanish \emph{blanco}).
\end{cousinweb}

\begin{grammarlens}[title={What you've built}]
One more word for telling, explaining and describing at B1.
\end{grammarlens}

\subsection*{Your turn}
\begin{itemize}
\raggedright
  \item \emph{Say it:} \emph{imbiancare}
  \item \emph{Say it:} \emph{imbiancare}, once more
  \item \emph{Say it:} a sentence of your own with \emph{imbiancare}
  \item \emph{Recall:} say the Italian for screwdriver, then the Italian for to fix, then say \emph{imbiancare} again
\end{itemize}

\begin{tcolorbox}[breakable,colback=teal!4,colframe=teal!35!black,title={Before you move on}]
What does imbiancare mean? (\textquotedblleft{}To paint (walls), to whitewash\textquotedblright{}.) Say it once more.
\end{tcolorbox}
